\section*{第 \ref{ch:b1:expression-control} 章——基因表达的调控}

\begin{solution}{exo:b1:expression-control:1}
\emph{lacI}（\omterm{def:b1:expression-control:operon}{阻遏蛋白}\omterm{def:b1:genomes:gene}{基因}）、\omterm{def:b1:gene-expression:transcription}{启动子}、\omterm{def:b1:expression-control:operon}{操纵基因}、\emph{lacZ}、
\emph{lacY}、\emph{lacA} 和 CAP 位点。无乳糖：关闭（\omterm{def:b1:expression-control:operon}{阻遏蛋白}
结合于\omterm{def:b1:expression-control:operon}{操纵基因}）。有乳糖、有葡萄糖：表达很弱（\omterm{def:b1:expression-control:operon}{阻遏蛋白}脱离，但无
cAMP--CAP）。有乳糖、无葡萄糖：完全表达。两种糖都没有：关闭。
\end{solution}

\begin{solution}{exo:b1:expression-control:2}
以反式作用的产物是一种可扩散的分子，能作用于任一
\omterm{def:b1:cell-unit-of-life:cell}{细胞}中其靶标的拷贝，即\omterm{def:b1:expression-control:operon}{阻遏蛋白}。以顺式作用的位点是
仅影响与其物理相连\omterm{def:b1:genomes:gene}{基因}的一段 \omterm{def:b1:nucleic-acids:chain}{DNA} 序列：
即\omterm{def:b1:expression-control:operon}{操纵基因}。
\end{solution}

\begin{solution}{exo:b1:expression-control:3}
延迟期约 \qty{36}{min}（从 \qty{2}{h} 到 \qty{2.6}{h}）；倍增
时间在葡萄糖上为 \qty{20}{min}，在乳糖上为 \qty{30}{min}。
\end{solution}

\begin{solution}{exo:b1:expression-control:4}
\omterm{def:b1:expression-control:enhancer}{增强子}：一段调控序列，常位于较远处，因子与之结合时可提高该\omterm{def:b1:genomes:gene}{基因}的
转录水平。\omterm{def:b1:expression-control:enhancer}{转录因子}：一种\omterm{def:b1:proteins:peptide}{蛋白质}，
具有 \omterm{def:b1:nucleic-acids:chain}{DNA} 结合结构域和激活或抑制结构域，可调控
转录。管家\omterm{def:b1:genomes:gene}{基因}：在每个\omterm{def:b1:cell-unit-of-life:cell}{细胞}中都以稳定水平表达的
\omterm{def:b1:genomes:gene}{基因}。
\end{solution}

\begin{solution}{exo:b1:expression-control:5}
$I^-$：组成型。$O^c$：组成型。$I^s$：不可诱导。
$I^-/I^+$：可诱导（$I^+$ \omterm{def:b1:expression-control:operon}{阻遏蛋白}以反式作用）。$I^s/I^+$：
不可诱导（$I^s$ 为显性）。$I^+\,O^c\,Z^- / I^+\,O^+\,Z^+$：可诱导
——唯一有功能的 $Z$ 位于正常的\omterm{def:b1:expression-control:operon}{操纵基因}之后。
\end{solution}

\begin{solution}{exo:b1:expression-control:6}
lac \omterm{def:b1:expression-control:operon}{操纵子}：糖（诱导物）与\omterm{def:b1:expression-control:operon}{阻遏蛋白}结合，使其
\emph{脱离}\omterm{def:b1:expression-control:operon}{操纵基因}——酶在其\omterm{def:b1:enzymes:enzyme}{底物}
存在时才被合成。色氨酸\omterm{def:b1:expression-control:operon}{操纵子}：\omterm{def:b1:proteins:aminoacid}{氨基酸}（辅阻遏物）与\omterm{def:b1:expression-control:operon}{阻遏蛋白}结合，
并\emph{使其能够}结合于\omterm{def:b1:expression-control:operon}{操纵基因}——酶在其产物
缺乏时才被合成。分解代谢途径由其输入物开启，
合成代谢途径则由其输出物关闭。
\end{solution}

\begin{solution}{exo:b1:expression-control:7}
在葡萄糖上：正常生长（葡萄糖不需要 CAP）。仅有乳糖时：
几乎不生长——无 cAMP 时，CAP 无法激活 lac
\omterm{def:b1:gene-expression:transcription}{启动子}，因此酶的量只有诱导水平的百分之几，即使
\omterm{def:b1:expression-control:operon}{阻遏蛋白}已脱离。两种糖同时存在时：先在葡萄糖上生长，随后停止，
不出现第二阶段。各种情况下酶水平都很低。
\end{solution}

\begin{solution}{exo:b1:expression-control:8}
衰减作用要求\omterm{def:b1:gene-expression:ribosome}{核糖体}在翻译信使 \omterm{def:b1:nucleic-acids:chain}{RNA} 的同时，
聚合酶仍在转录它，使\omterm{def:b1:gene-expression:ribosome}{核糖体}的位置
能够塑造聚合酶身后的 \omterm{def:b1:nucleic-acids:chain}{RNA}。在真核生物中，
转录发生在\omterm{def:b1:eukaryotic-cell:nucleus}{细胞核}内，翻译发生在\omterm{def:b1:eukaryotic-cell:organelle}{细胞溶胶}中，
mRNA 在加工并
输出之前不会被翻译。
\end{solution}

\begin{solution}{exo:b1:expression-control:9}
\omterm{def:b1:expression-control:enhancer}{增强子}作用距离较远，无论朝向如何、位于哪一侧都可起作用，
因而不能靠被读取或直接定位聚合酶发挥作用；
它们通过结合因子起作用，这些因子借助 \omterm{def:b1:nucleic-acids:chain}{DNA} 环
与\omterm{def:b1:gene-expression:transcription}{启动子}接触，距离与朝向对此
影响很小。
\end{solution}

\begin{solution}{exo:b1:expression-control:10}
$5000\times 4096\times 4 = \num{8.2e7}$ 个 \omterm{def:b1:water-small-molecules:atp}{ATP}，占
预算的 \qty{0.8}{\%}。生长速率降低 \qty{0.8}{\%}：每代中突变体的
份额按因子 $2^{-0.008} = 0.9945$ 下降。从 $1:1$ 变至
$1:99$：$0.9945^n = 1/99$，$n = \ln 99/0.00555 = 830$ 代——
连续培养数周。
\end{solution}

\begin{solution}{exo:b1:expression-control:11}
每种\omterm{def:b1:cell-unit-of-life:cell}{细胞}类型都带有一组特定的\omterm{def:b1:expression-control:enhancer}{转录因子}，它们
与其\omterm{def:b1:genomes:gene}{基因}的\omterm{def:b1:expression-control:enhancer}{增强子}结合，并招募乙酰化酶，使其
\omterm{def:b1:genomes:nucleosome}{染色质}开放，而其他类型的\omterm{def:b1:genomes:gene}{基因}则被甲基化，
并被包装为任何因子都无法触及的异\omterm{def:b1:genomes:nucleosome}{染色质}。这些因子
相互维持彼此的表达（一个具有稳定状态的网络），并且
\omterm{def:b1:genomes:nucleosome}{染色质}标记在复制时一并传递，因此子\omterm{def:b1:cell-unit-of-life:cell}{细胞}同时继承
因子以及开放与关闭的区域。\omterm{def:b1:genomes:genome}{基因组}是
相同的；可及的\omterm{def:b1:genomes:genome}{基因组}却不然，这种不可及性
可遗传。
\end{solution}

\begin{solution}{exo:b1:expression-control:12}
细菌的调控与其所处\omterm{def:b1:organism-environment:organism}{环境}相适应：糖、\omterm{def:b1:proteins:aminoacid}{氨基
酸}、温度等信号与\omterm{def:b1:expression-control:operon}{阻遏蛋白}或 sigma 因子结合，应答
在数分钟内完成，信号一旦消失便随即逆转——
表达随\omterm{def:b1:organism-environment:organism}{环境}而变化。多\omterm{def:b1:cell-unit-of-life:cell}{细胞}\omterm{def:b1:organism-environment:organism}{生物体}中的\omterm{def:b1:cell-unit-of-life:cell}{细胞}也会对
信号作出应答，但其主要决定是在发育过程中作出的，并
被锁定于\omterm{def:b1:genomes:nucleosome}{染色质}之中：其能表达的内容由所继承的因子
与 \omterm{def:b1:nucleic-acids:chain}{DNA} 上的标记决定，且仅能通过分裂而改变，
大多不可逆。细菌是当下\omterm{def:b1:organism-environment:organism}{环境}的读取者；而分化\omterm{def:b1:cell-unit-of-life:cell}{细胞}
则具有一种记忆，这种记忆即其身份。
\end{solution}

\begin{solution}{pb:b1:expression-control:1}
\textbf{1.} $0.5\times 10^{-3}\times 180 = \qty{0.09}{g/L} =
\qty{9e-5}{g/mL}$.
\textbf{2.} $9\times 10^{-5}/1.5\times 10^{-12} = \num{6e7}$ 次分裂。
\textbf{3.} $10^7 + \num{6e7} = \num{7e7}$ 个\omterm{def:b1:cell-unit-of-life:cell}{细胞}/毫升。
\textbf{4.} $7 = 2^n$：$n = 2.8$ 次倍增，\qty{56}{min}。
\textbf{5.} 关闭：乳糖（别乳糖）已使\omterm{def:b1:expression-control:operon}{阻遏蛋白}脱离，
但葡萄糖使 cAMP 维持在较低水平，因此 CAP 未结合，\omterm{def:b1:gene-expression:transcription}{启动子}
几乎沉默；同时通透酶稀少，进入\omterm{def:b1:cell-unit-of-life:cell}{细胞}的乳糖很少。
\textbf{6.} cAMP 升高并与 CAP 结合，后者再与\omterm{def:b1:gene-expression:transcription}{启动子}结合。
\textbf{7.} 负调控已解除，但正调控缺失：
没有 CAP 的弱\omterm{def:b1:gene-expression:transcription}{启动子}很少起始转录（\omterm{def:b1:expression-control:cap}{分解代谢物
阻遏}）。
\textbf{8.} $5000/1800 = 2.8$ 个四聚体/秒；$2.8\times 4096 =
\num{11400}$ 个残基/秒。
\textbf{9.} 容量 $20\,000\times 20 = \num{4e5}$ 个残基/秒：
占\omterm{def:b1:gene-expression:ribosome}{核糖体}的 \qty{3}{\%}。
\textbf{10.} $5000\times 4096\times 110\times 1.66\times 10^{-24} =
\qty{3.7e-15}{g}$：占\omterm{def:b1:cell-unit-of-life:cell}{细胞}\omterm{def:b1:proteins:peptide}{蛋白质}的 \qty{2.5}{\%}。
\textbf{11.} $5000/5 = 1000$.
\textbf{12.} 没有通透酶，乳糖无法穿过膜；
\emph{lacY}$^-$ 突变体虽有通透酶，但\omterm{def:b1:cell-unit-of-life:cell}{细胞}内没有\omterm{def:b1:enzymes:enzyme}{底物}，
不能在乳糖上生长（诱导也不良，因为诱导物是在
\omterm{def:b1:cell-unit-of-life:cell}{细胞}内由乳糖生成的）。
\textbf{13.} $1.0\times 10^{-3}\times 342 = \qty{0.342}{g/L}$，相当于
每摩尔 $2\times 180 = \qty{360}{g}$ 葡萄糖，即 \qty{0.36}{g/L}
$= \qty{3.6e-4}{g/mL}$.
\textbf{14.} $3.6\times 10^{-4}/1.5\times 10^{-12} = \num{2.4e8}$
次分裂；最终为 $\num{7e7} + \num{2.4e8} = \num{3.1e8}$ 个\omterm{def:b1:cell-unit-of-life:cell}{细胞}/毫升。
\textbf{15.} $3.1/0.7 = 4.4 = 2^n$：$n = 2.1$ 次倍增，\qty{64}{min}。
\textbf{16.} 葡萄糖 \qty{56}{min}，延迟期 \qty{30}{min}，乳糖
\qty{64}{min}：共计 \qty{150}{min}。
\textbf{17.} 乳糖必须先被水解，这一额外的酶促步骤
其速率限制了葡萄糖的供应；而通透酶利用
质子梯度将其输入\omterm{def:b1:cell-unit-of-life:cell}{细胞}，这是葡萄糖转运所不需付出的代价。
\textbf{18.} $I^-$：酶的水平始终为 \num{5000}（组成型，一旦
葡萄糖耗尽即如此；在葡萄糖上 CAP 仍限制其表达，因此水平
处于中间值）；其延迟期较短或不存在，因此曲线上的
停顿较小；在葡萄糖上，这一代价使其生长略慢。
\textbf{19.} CAP$^-$：在葡萄糖上生长至 $\num{7e7}$，然后停止：
lac \omterm{def:b1:expression-control:operon}{操纵子}无法被激活，酶的数量保持在 5 个分子左右，并且
乳糖从不被利用。
\textbf{20.} 仅有乳糖时：诱导后两者的酶水平并无差异（两者
都完全开启），生长上也基本无差异；$O^c$ 只是缺少
诱导延迟。仅有葡萄糖时：$O^c$ 无用地合成该酶
（在 CAP 所允许的范围内），野生型则不然。
\textbf{21.} $5000\times 4096\times 4 = \num{8.2e7}$ 个 \omterm{def:b1:water-small-molecules:atp}{ATP}：\qty{0.8}{\%}
占 $10^{10}$ 的比例。
\textbf{22.} 世代时间延长 \qty{0.8}{\%}：生长速率
降低 \qty{0.8}{\%}，$\delta = 0.008$ 次倍增/代。
\textbf{23.} $r = 2^{-0.008} = 0.99447$; $r^n = 0.111$: $n =
\ln 0.111/\ln 0.99447 = 2.20/0.00555 = 400$ 代。
\textbf{24.} 在乳糖上，酶本来就是需要的，因此该突变体
不付出代价，并且省去了延迟期：它不受选择淘汰，
有时甚至受到正选择；而实验室菌株的培养
是出于方便，而非竞争。
\textbf{25.} 诱导倍数为 1000（从 5 到 \num{5000} 个分子）；
不必要的合成耗费 \qty{0.8}{\%} 的 \omterm{def:b1:water-small-molecules:atp}{ATP} 预算；
实验约持续两个半小时。
\end{solution}
